\documentclass[a4,11pt]{aleph-notas}
% Se puede ver la documentación aquí:
% https://github.com/alephsub0/LaTeX_aleph-notas

% -- Paquetes adicionales
\usepackage{enumitem}
\usepackage{url}
\usepackage{array}
\usepackage{booktabs}
\usepackage{longtable}
\usepackage{aleph-comandos}
\hypersetup{
    urlcolor=blue,
    linkcolor=blue,
}



% -- Datos
\institucion{Facultad de Ciencias Exactas, Naturales y Ambientale}
\carrera{Catálogo STEM}
\asignatura{Matemática III}
\tema{Plan tema 01: Funciones vectoriales y curvas}
\autor{Andrés Merino}
\fecha{Periodo 2026-1}

\logouno[0.16\textwidth]{Logos/logoPUCE_04_ac}
\definecolor{colortext}{HTML}{1957d1}
\definecolor{colordef}{HTML}{1957d1}
\fuente{montserrat}


\begin{document}

\encabezado

%%%%%%%%%%%%%%%%%%%%%%%%%%%%%%%%%%%%%%%%
\section*{Resultado de Aprendizaje}
%%%%%%%%%%%%%%%%%%%%%%%%%%%%%%%%%%%%%%%%

%%%%%%%%%%%%%%%%%%%%%%%%%%%%%%%%%%%%%%%%
\subsection*{RdA de la asignatura:}
\begin{itemize}[leftmargin=*]
    \item \textbf{RdA 1:} Identificar los conceptos, teoremas y propiedades de funciones vectoriales, derivadas parciales e integral vectorial.
    % \item \textbf{RdA 2:} Calcular derivadas parciales, integrales de funciones vectoriales con el apoyo de asistentes computacionales.
    % \item \textbf{RdA 3:} Aplicar distintos tópicos del Cálculo Vectorial para la solución de problemas reales en diferentes contextos.
\end{itemize}

%%%%%%%%%%%%%%%%%%%%%%%%%%%%%%%%%%%%%%%%
\subsection*{RdA del tema:}
\begin{itemize}[leftmargin=*]
    \item Clasificar funciones vectoriales en trayectorias, campos escalares y campos vectoriales según sus dimensiones de dominio e imagen.
    \item Describir e interpretar los conjuntos de nivel de un campo escalar.
    \item Calcular límites, derivadas e integrales de trayectorias componente a componente.
    \item Determinar el vector tangente y la recta tangente a una curva parametrizada.
    \item Interpretar la posición, velocidad, rapidez y aceleración de una partícula como una trayectoria y sus derivadas.
\end{itemize}

%%%%%%%%%%%%%%%%%%%%%%%%%%%%%%%%%%%%%%%%
\section*{Introducción}
%%%%%%%%%%%%%%%%%%%%%%%%%%%%%%%%%%%%%%%%

\paragraph{Pregunta inicial:}
¿Cómo se puede describir matemáticamente la trayectoria de un objeto que se mueve en el espacio, y qué información geométrica se puede extraer de esa descripción?

%%%%%%%%%%%%%%%%%%%%%%%%%%%%%%%%%%%%%%%%
\section*{Desarrollo}
%%%%%%%%%%%%%%%%%%%%%%%%%%%%%%%%%%%%%%%%

%%%%%%%%%%%%%%%%%%%%%%%%%%%%%%%%%%%%%%%%
\subsection*{Actividad 1: Funciones vectoriales y conjuntos de nivel}
La clase inicia con ejemplos concretos de funciones de $\R^n$ en $\R^m$: temperatura en un punto del espacio (campo escalar), posición de una partícula (trayectoria) y un campo de vientos (campo vectorial). A través de una clase magistral se trabajan las definiciones de función vectorial, componentes, conjuntos de nivel y operaciones entre funciones.

\paragraph{¿Cómo lo haremos?}
\begin{itemize}[leftmargin=*]
    \item \textbf{Motivación:} se presentan ejemplos físicos de cada tipo de función vectorial y se discute cómo representarlos matemáticamente.
    \item \textbf{Clase magistral:} se definen trayectorias, campos escalares y vectoriales, componentes, conjuntos de nivel (curvas y superficies de nivel) y operaciones. Se usa el resumen \href{https://andres-merino.github.io/MatematicaIII-05-N0051/01Resumen/Resumen01.pdf}{Resumen01.pdf}.
    \item \textbf{Resolución de ejercicios:} los estudiantes describirán e interpretarán conjuntos de nivel de campos escalares, usando el documento \href{https://andres-merino.github.io/MatematicaIII-05-N0051/04Ejercicios/Ejercicios01.pdf}{Ejercicios01.pdf} (ejercicios 1--6).
\end{itemize}

\paragraph{Verificación de aprendizaje:}
\begin{itemize}[leftmargin=*]
    \item ¿Qué diferencia hay entre una trayectoria y un campo escalar?
    \item ¿Qué representa geométricamente el conjunto de nivel $L_c(f)$ cuando $n=2$?
    \item Dado $f(x,y)=x^2+y^2$, ¿cómo son los conjuntos de nivel de $f$?
\end{itemize}

%%%%%%%%%%%%%%%%%%%%%%%%%%%%%%%%%%%%%%%%
\subsection*{Actividad 2: Cálculo en trayectorias y curvas}
Se trabaja el cálculo diferencial e integral de trayectorias, haciendo énfasis en que las propiedades se reducen al cálculo componente a componente. Se introduce la noción de curva como imagen de una trayectoria continua y se estudia el vector tangente y las parametrizaciones equivalentes.

\paragraph{¿Cómo lo haremos?}
\begin{itemize}[leftmargin=*]
    \item \textbf{Clase magistral:} se definen límite, continuidad, derivada e integral de trayectorias, los teoremas de derivación (producto escalar, vectorial, regla de la cadena) y los Teoremas Fundamentales del Cálculo. Se usa el resumen \href{https://andres-merino.github.io/MatematicaIII-05-N0051/01Resumen/Resumen01.pdf}{Resumen01.pdf}.
    \item \textbf{Resolución de ejercicios:} los estudiantes calcularán límites, derivadas e integrales de trayectorias y determinarán rectas tangentes a curvas, usando \href{https://andres-merino.github.io/MatematicaIII-05-N0051/04Ejercicios/Ejercicios01.pdf}{Ejercicios01.pdf} (ejercicios 7--21).
\end{itemize}

\paragraph{Verificación de aprendizaje:}
\begin{itemize}[leftmargin=*]
    \item ¿Cómo se calcula $\lim_{t\to a}\alpha(t)$ cuando $\alpha:\R\to\R^3$?
    \item Si $\|\alpha(t)\|=k$ para todo $t$, ¿qué relación hay entre $\alpha(t)$ y $\alpha'(t)$?
    \item ¿Qué condición debe cumplir $\alpha'(t)$ para que exista recta tangente en $\alpha(t)$?
\end{itemize}

%%%%%%%%%%%%%%%%%%%%%%%%%%%%%%%%%%%%%%%%
\subsection*{Actividad 3: Movimiento en el espacio}
A partir del concepto de trayectoria diferenciable, se introduce la interpretación física de sus derivadas: la primera como velocidad (vector) y rapidez (escalar), y la segunda como aceleración. Se presentan los tipos de movimiento —rectilíneo, circular y planar— como casos especiales de la forma de la función posición.

\paragraph{¿Cómo lo haremos?}
\begin{itemize}[leftmargin=*]
    \item \textbf{Motivación:} se presenta el movimiento de un proyectil o de un planeta y se discute qué información matemática se puede extraer de su trayectoria.
    \item \textbf{Clase magistral:} se definen posición, velocidad, rapidez y aceleración de una trayectoria; se caracterizan los movimientos rectilíneo, circular y planar. Se usa el resumen \href{https://andres-merino.github.io/MatematicaIII-05-N0051/01Resumen/Resumen02.pdf}{Resumen02.pdf}.
    \item \textbf{Resolución de ejercicios:} los estudiantes calcularán vectores velocidad y aceleración, e identificarán el tipo de movimiento, usando el documento \href{https://andres-merino.github.io/MatematicaIII-05-N0051/04Ejercicios/Ejercicios02.pdf}{Ejercicios02.pdf}.
    \item \textbf{Implementación computacional:} visualización de trayectorias con sus vectores velocidad y aceleración mediante el cuaderno \href{https://andres-merino.github.io/MatematicaIII-05-N0051/02ResumenPython/ResumenPy02.pdf}{ResumenPy02.pdf}.
\end{itemize}

\paragraph{Verificación de aprendizaje:}
\begin{itemize}[leftmargin=*]
    \item Dada $\alpha(t)=(\cos t, \sen t, t)$, ¿cuáles son $v(t)$, $\rho(t)$ y $a(t)$?
    \item ¿Qué forma debe tener $\alpha(t)$ para que el movimiento sea rectilíneo?
    \item ¿Puede la rapidez de una partícula ser constante aunque su aceleración sea no nula?
\end{itemize}

%%%%%%%%%%%%%%%%%%%%%%%%%%%%%%%%%%%%%%%%
\section*{Cierre}
%%%%%%%%%%%%%%%%%%%%%%%%%%%%%%%%%%%%%%%%

\paragraph{Pregunta de investigación:}
\begin{enumerate}[leftmargin=*]
    \item ¿Cómo se usa la parametrización de curvas en animación por computadora o en robótica para describir el movimiento de un brazo mecánico?
    \item ¿Qué relación existe entre los conjuntos de nivel de una función y las líneas de contorno en un mapa topográfico?
\end{enumerate}

\paragraph{Para el próximo tema:}
 Ver el video:
\begin{itemize}[leftmargin=*]
    \item \href{https://youtu.be/VIqA8U9ozIA?si=G4a9bhE6ieUn622f}{Torsion: How curves twist in space, and the TNB or Frenet Frame}
\end{itemize}

\end{document}
